\section{Introduction}

Speculative decoding accelerates large language model (LLM) inference by
letting a small draft model propose multiple tokens that are then verified by a
larger target model in parallel \citep{leviathan2023speculative}. The resulting
speedup is lossless at the output layer, but it is not free: it depends
critically on the degree of alignment between the draft distribution and the
target verifier. When draft-target agreement is high, long blocks of drafted
tokens are accepted and throughput improves substantially. When agreement
collapses, the verifier commits fewer drafted tokens, target-side correction
work increases, and the latency benefits of speculative decoding shrink.

Most prior work treats this acceptance bottleneck as an optimization problem to
be fixed. Recent systems improve the draft, adapt draft length, redesign
verification, or otherwise try to raise acceptance rates
\citep{liu2023online,goel2024direct,agrawal2024adaedl,bachmann2025judge,wang2025alignment,zhou2026hsd}.
This paper studies the complementary adversarial question: can an attacker
exploit the same acceptance bottleneck from the input side? In particular, can
a short prompt prefix preserve user-visible task behavior while forcing the
verifier into a substantially less efficient regime?

This attack surface is natural in deployment. The attacker does not need model
weights, scheduler privileges, or a custom decoding stack at inference time.
Instead, the attacker controls only the prompt text. Because speculative
decoding ties latency directly to the relation between the draft distribution
$q$ and the target distribution $p$, prompt-side inputs can act as
systems-level levers: they can induce overconfident draft proposals that the
target model repeatedly rejects, thereby converting speculative acceleration
into additional verifier work.

The technical obstacle is discrete. Gradient-based prompt optimization operates
most naturally over continuous token relaxations, but deployment requires a
discrete projected prefix. The main question is therefore not simply whether a
smooth surrogate can be written down, but when optimizing that surrogate yields
a strong \emph{discrete} attack after projection. This continuous-to-discrete
gap is the central methodological challenge for prompt-side attacks on
speculative decoding.

Our starting point is the token-wise verifier acceptance rule
\[
\alpha_i = \min\!\left(1, \frac{p(y_i \mid h_i)}{q(y_i \mid h_i)}\right),
\]
where $h_i$ is the speculative history and $y_i$ is the drafted token at step
$i$. Under draft sampling, the one-step expected acceptance satisfies
\[
\mathbb{E}[\alpha_i] = 1 - \tv(p_i, q_i),
\]
which immediately motivates an adversarial objective: maximize early
draft-target disagreement, since early verifier rejection is the most damaging
to end-to-end speculative efficiency. We formalize this intuition with a
continuous objective over relaxed prefixes, show that unconditional projection
guarantees are impossible via a prompt-waywardness proposition, and prove a
relaxation-gap bound under a TV-Lipschitz assumption that directly motivates a
projection-distortion penalty in the implemented optimizer.

On the algorithmic side, we instantiate this theory in two stages. The first
stage is a verifier-aligned soft-prefix optimizer that targets first-token
collapse at the assistant boundary using a straight-through relaxed prefix, a
soft-collapse term, TV, reverse KL, entropy sharpening, and an explicit
projection-distortion penalty. The second stage strengthens the attack into a
calibration-batch universal adversarial prefix optimized over target-generated
multi-timestamp rollouts. This universal-prefix pipeline combines a weighted
trajectory-level objective with a gradient-guided discrete seed-search stage,
allowing the optimizer to discover stronger discrete initializations before
continuous refinement.

We evaluate the attack in the reproduced HSD workspace using
\texttt{Qwen2.5-0.5B-Instruct-GPTQ-Int8} as the draft model and
\texttt{Qwen2.5-14B-Instruct-GPTQ-Int8} as the target model. The initial
Palmetto soft-prefix search consistently recovers the same five-token discrete
boundary attack across a sweep of projection-penalty strengths, reducing exact
first-token acceptance from $\alpha_1 = 0.1597$ for the manual comparison seed
to $\alpha_1 = 0.0288$ on the optimized search context. On a matched
263-sample token-wise GSM8K benchmark, that boundary attack raises mean sample
time from $23.10$s to $29.81$s while reducing aggregate acceptance from
$0.597$ to $0.509$ and block efficiency from $5.832$ to $4.961$.

The stronger result comes from the calibration-batch universal-prefix attack.
The best accuracy-preserving universal prefix raises mean sample time to
$38.61$s, lowers aggregate acceptance to $0.390$, lowers block efficiency to
$3.731$, and reduces decoding speed to $51.68$, while matching the benign
token-wise GSM8K accuracy exactly at $222/263$. In other words, a short
prompt-side prefix can preserve task accuracy while extracting a large
inference-cost penalty from a lossless speculative-decoding stack.

Our contributions are:
\begin{itemize}
\item We formulate prompt-side speculative-decoding slowdown as a
continuous-to-discrete adversarial optimization problem over verifier
acceptance rather than output semantics.
\item We prove a prompt-waywardness impossibility result and a relaxation-gap
bound that makes the projection distortion term in the optimizer theoretically
explicit.
\item We implement a verifier-aligned soft-prefix attack and a
calibration-batch universal-prefix extension with gradient-guided discrete
seed search, both of which are faithful to the released code path.
\item We show on a reproduced HSD/Qwen stack that prompt-only attacks can
materially degrade speculative efficiency, and that the strongest
accuracy-preserving universal prefix raises end-to-end runtime by $67.1\%$ on a
matched token-wise benchmark.
\end{itemize}

The broader message is that speculative decoding exposes a distinct
prompt-controlled systems vulnerability. Lossless output semantics do not imply
cost robustness. A user-controlled prefix can leave the final answer intact
while forcing the verifier to do substantially more work.
